% SPDX-License-Identifier: Apache-2.0
%
% TxHDL, step by step: a tutorial for somebody meeting the language for
% the first time. Built by //docs:tutorial. Every listing is included
% from the tree, and every printout and figure is what the build made
% by running the example, so the tutorial cannot say what the code
% does not do. Two columns, as the paper is; the listings are set at
% the paper's size, and the waveforms are drawn at column width.
\documentclass[journal,9pt]{IEEEtran}

% Listings of Rust set by rustfmt at 80 columns: 5.6pt is what fits
% the frame of a column (issue 1061).
\newcommand{\listingsize}{\fontsize{5.6}{6.6}\selectfont}
% SPDX-License-Identifier: Apache-2.0
%
% The house preamble, shared by every document in this directory. Loaded
% after \documentclass, before \title. Keeping it in one file is what
% stops two articles drifting apart in font, listing style or figure
% style.
% Computer Modern. IEEEtran selects Times (ptm) by default, so the face has
% to be chosen explicitly.
%
% Latin Modern is the Computer Modern variety used here, and the choice is
% forced rather than aesthetic. Under T1 encoding, plain `cmr` has no Type 1
% outlines unless cm-super is installed, and it is not in the pinned
% distribution, so pdflatex falls back to EC bitmaps and every glyph in the
% PDF becomes a Type 3 font. That was measured: the first build produced a
% document with 10 Type 3 fonts and no Type 1. Latin Modern is Computer
% Modern redrawn with a full T1 repertoire and real .pfb outlines, so the
% same design comes out scalable.
\usepackage[T1]{fontenc}
\usepackage[utf8]{inputenc}
\usepackage{lmodern}
% microtype is not in the pinned TeX distribution. emergencystretch alone
% keeps the two-column measure from producing overfull lines.
\setlength{\emergencystretch}{3em}

\usepackage{amsmath}
\usepackage{amssymb}
\usepackage{amsthm}
\usepackage{booktabs}
\usepackage{array}
% enumitem is not in the pinned TeX distribution; the list spacing below
% does what \setlist would have done.
\usepackage{float}
\usepackage{xcolor}
\usepackage{listings}
\usepackage{tikz}
\usetikzlibrary{arrows.meta,positioning,fit,backgrounds,calc}
\usepackage[hidelinks,bookmarks=true]{hyperref}

% Numbered, definition-style box for a rule the reader is meant to apply.
\theoremstyle{definition}
\newtheorem{rulebox}{Rule}
\newtheorem{example}{Example}

% Unnumbered asides.
\theoremstyle{remark}
\newtheorem*{tip}{Tip}
\newtheorem*{pitfall}{Pitfall}

% Tighter lists, in place of enumitem's \setlist.
\setlength{\topsep}{2pt}
\setlength{\itemsep}{1pt}
\setlength{\parsep}{0pt}

% Code in the running text. The typewriter face under T1 ligates `<<`
% and `>>` into guillemets, so a shift printed as \code{<<} came out as
% a French quotation mark (issue 571). microtype's \DisableLigatures is
% not in the pinned distribution, but the primitive it rests on is
% pdfTeX's own: \pdfnoligatures strips every ligature from the font it
% is given, and the current font inside \texttt is the typewriter face
% at the size in use, so each size is stripped the first time \code
% sets it. Code wants no ligature at all: two quotes are two quotes.
\newcommand{\code}[1]{\texttt{\ifdefined\pdfnoligatures\pdfnoligatures\font\fi #1}}
\newcommand{\term}[1]{\emph{#1}}

% Syntax highlighting. Four roles, distinguishable in colour and still
% distinguishable in greyscale, because an IEEE article gets printed: the
% keyword is the only bold role, the comment is the only italic one, and the
% two remaining roles differ in lightness as well as in hue.
\definecolor{kwblue}{RGB}{20,60,150}
\definecolor{cmtgreen}{RGB}{90,120,90}
\definecolor{strred}{RGB}{150,50,40}
\definecolor{numgray}{gray}{0.45}
\definecolor{framegray}{gray}{0.70}
\definecolor{bgshade}{gray}{0.97}
% A darker fill for figure cells. The listing background has to stay very
% light to keep code readable; a figure cell has to be clearly shaded or the
% distinction it draws is invisible in print.
\definecolor{cellshade}{gray}{0.90}

% The listing size is the document's choice, because it depends on the
% column: a two-column article sets `\listingsize` to 6pt and keeps its
% listings within 70 columns, or to 5.6pt when it includes source that
% rustfmt sets at 80, which is what fits a column's frame (issue 1061);
% a one-column document keeps the body size and includes source files
% at 80. `//docs:frame_test` checks every listing line against its frame. Lines are never broken; a line that does not fit
% overflows the frame, which is visible, rather than wrapping, which is
% not.
\providecommand{\listingsize}{\scriptsize}

% `'` is deliberately not a string delimiter: the language uses it for loop
% labels, as Rust does, so treating it as a quote would swallow the rest of
% the line.
\lstdefinestyle{txhdl}{
  basicstyle=\ttfamily\listingsize,
  keywordstyle=\bfseries\color{kwblue},
  commentstyle=\itshape\color{cmtgreen},
  stringstyle=\color{strred},
  numberstyle=\tiny\color{numgray},
  morekeywords={mod,use,pub,const,type,struct,enum,trait,impl,fn,async,let,mut,
    match,if,else,loop,while,for,in,break,continue,return,as,await,
    module,bus,channel,transaction,protocol,pipeline,flow,seq,config,
    port,stage,pipe,instance,connect,bind,tag,domain,blackbox,foreign,
    property,role,signal,handler,true,false,
    sequence,interface,modport,implements,package,import,var,wire,func,proc,
    case,default,next,fork,branch,join,state,fsm},
  morecomment=[l]{//},
  morestring=[b]",
  showstringspaces=false,
  columns=fullflexible,
  keepspaces=true,
  breaklines=false,
  % Framed, so a listing is bounded rather than floating in the column.
  frame=single,
  rulecolor=\color{framegray},
  framesep=3pt,
  backgroundcolor=\color{bgshade},
  xleftmargin=3pt,
  xrightmargin=3pt,
  aboveskip=6pt,
  belowskip=6pt,
  % Every listing is numbered and captioned, so it can be referred to by
  % number. `caption` without `float` numbers the listing and leaves it
  % where it was written, which is what a listing in running prose needs.
  captionpos=b,
  abovecaptionskip=2pt,
  belowcaptionskip=6pt,
}

% Figures drawn as text. Framed the same way, and with the highlighting
% switched off, because a box-and-arrow diagram is not code and half its
% words turning blue reads as an accident. Setting a style adds keys rather
% than clearing the ones already set, so the three roles are neutralised by
% hand rather than by leaving them out. Program output is printed in this
% style too, and a netlist's assignment can run past the frame, so a line
% that does not fit is broken and the rest indented; source never is.
\lstdefinestyle{ascii}{
  basicstyle=\ttfamily\listingsize,
  keywordstyle=\ttfamily,
  commentstyle=\ttfamily,
  stringstyle=\ttfamily,
  columns=fullflexible,
  keepspaces=true,
  breaklines=true,
  breakindent=2em,
  frame=single,
  rulecolor=\color{framegray},
  framesep=3pt,
  backgroundcolor=\color{bgshade},
  xleftmargin=3pt,
  xrightmargin=3pt,
  aboveskip=6pt,
  belowskip=6pt,
}
\lstset{style=txhdl}

% House TikZ styles. `shorten` lives on the arrow styles rather than on each
% arrow, so every arrow in the document gets the gap, including ones added
% later. TikZ clips a path at a node's border, so without it an arrowhead
% butts against the box with no gap at all. Keep `node distance` well above
% twice the shorten, or the arrow shrinks to nothing.
\tikzset{
  box/.style={draw,rounded corners=2pt,align=center,inner sep=4pt,
              font=\footnotesize},
  boxf/.style={box,fill=bgshade},
  lbl/.style={font=\scriptsize\itshape,align=center},
  fl/.style={-{Stealth[length=4pt]},shorten >=2.5pt,shorten <=2.5pt},
  fld/.style={fl,dashed},
}


% The contents list: the class gives a section's number 2.75em, which
% a Roman numeral past thirty overruns onto the title, so the box is
% wider here, and a subsection's entry starts where a title does.
\makeatletter
\def\l@section#1#2{\addpenalty{\@secpenalty}\addvspace{1.0em plus 1pt}%
    \@tempdima 5.75em \begingroup \parindent \z@ \rightskip \@pnumwidth%
    \parfillskip-\@pnumwidth {\bfseries\leavevmode #1}\hfil\hbox to\@pnumwidth{\hss #2}\par%
    \endgroup}
\def\l@subsection{\@dottedtocline{2}{5.75em}{5em}}
\def\l@subsubsection{\@dottedtocline{3}{10.75em}{5.5em}}
\makeatother


\InputIfFileExists{buildstamp.tex}{}{}
\providecommand{\buildstamp}{Draft build (outside the Bazel flow).}

\title{TxHDL, Step by Step}

\author{Filip Filmar and Dragiša Janković%
\thanks{Both authors are independent. Filip Filmar is the
corresponding author, at f@hdlfactory.com; Dragiša Janković is
reached at linkedin.com/in/dragisa-jankovic.}%
\thanks{This tutorial guides developers with Rust familiarity through
authoring synthesizable hardware, running simulations, inspecting waveforms,
and deriving verified netlists. Every listing is extracted directly from the
repository \code{HDL/txhdl} during the build, and execution logs reflect actual
testbench output. Theoretical foundations appear in~\cite{embedding}, runtime
mechanics in~\cite{runtime}, reference examples in~\cite{examples}, and the
complete document corpus in~\cite{cover}.}%
\thanks{The authors used a large language model, Claude, as an
assistant in exploring concepts and in drafting and constructing
the documents and implementations; every repository commit documents
this collaboration and includes the prompt sequence verbatim.}%
\thanks{\buildstamp}}

\markboth{TxHDL, Step by Step}{Filmar and Janković: TxHDL, Step by Step}

\begin{document}
\maketitle

\begin{abstract}
TxHDL is a hardware description language implemented directly as a Rust
library. A hardware design is a struct whose sequential state resides in typed
registers, governed by an asynchronous execution function that synchronizes to
clock edges, reads sampled inputs, and asserts synchronous drives. A single
source file simulates natively in Rust, emits execution waveforms, and lowers to
cycle-equivalent Verilog and VHDL netlists verified automatically against
simulation traces. This tutorial introduces the system through eight executable
milestones: basic counters, discrete-event time, conditional multiplexing,
transactional channels, structured port bundles, multi-clock domain
hierarchies, waveform visualization, and netlist verification. A final
milestone demonstrates contributing new modules to the verified codebase.
\end{abstract}

\begin{IEEEkeywords}
Hardware description languages, Rust, tutorial, simulation, Verilog,
VHDL.
\end{IEEEkeywords}

\section{Before You Start}
\label{sec:t-start}

Toolchain prerequisites are minimal: \code{bazelisk}, which provisions the
pinned Bazel binary, and Git. Compilers, HDL simulators, formal model
checkers, and the \LaTeX{} engines that typeset this document are fetched
hermetically by Bazel. Clone the repository and execute the test suite:

\begin{lstlisting}[style=ascii]
bazel build //...
bazel test  //...
\end{lstlisting}

The initial invocation caches hermetic toolchains; subsequent builds complete in
seconds. Successful execution confirms a fully operational environment.

The core mental model rests on three principles. A \emph{unit} is a struct whose
fields declare sequential registers, animated by an \code{async fn run} process.
The process synchronizes to active clock edges, samples registered and input
values, and asserts combinational and registered drives: register updates take
effect at cycle boundaries while combinational drives propagate immediately.
When constrained to the lowerable language subset, a single unit source compiles
both as an executable Rust model and as a synthesizable Verilog and VHDL
netlist.

\section{Step One: A Counter}
\label{sec:t-counter}

Listing~\ref{lst:t-counter} presents a self-contained unit: a 3-bit counter that
advances while \code{enable} is asserted and emits a synchronous \code{tick}
pulse at terminal count seven.

\lstinputlisting[rangebeginprefix=//\ begin\{,rangebeginsuffix=\},rangeendprefix=//\ end\{,rangeendsuffix=\},includerangemarker=false,linerange=unit-unit,caption={The counter: a register, a wait, two reads, a predicated drive and an output. From \code{ex\_cheat.rs}.},label={lst:t-counter}]{lib/examples/ex_cheat.rs}

The struct definition encapsulates sequential state: a single 3-bit register,
\code{n}. \code{Reg<T>} wraps any value type; \code{U<3>} denotes an unsigned
3-bit integer. Deriving \code{Trace} exposes the register by name within
simulation waveforms, while \code{Default} specifies its reset state.

The \code{impl} block exposes port signatures: a scalar input \code{enable:
In<Bit>} and an output \code{tick: Out<Bit>}. \code{In} and \code{Out}
represent paired endpoints of a combinational wire. The \code{\#[lower]}
procedural attribute derives synthesizable netlists directly. Without this
attribute, units remain simulatable provided the header declares explicit port
interfaces (\code{impl Unit<In, Out> for Counter}).

The process body executes an infinite \code{loop} initiated by an edge wait:
\code{DefaultClock::rising().await}. Logic following this wait executes within a
single clock cycle. Input \code{enable.get()} reflects the wire state sampled at
the clock edge, while \code{self.n} reads the registered count. Macro
\code{with!} governs sequential updates: when \code{enable} is true, \code{n}
latches \code{n + 1} at the cycle boundary, using struct-literal syntax with
update operator \code{<=}. Combinational output \code{tick.set} asserts during
the active cycle based on the state sampled at edge arrival. This
structure---a wait, reads, combinational evaluations, and single-assignment
drives---defines the canonical shape of all lowerable units.

Execute the design via \code{bazel run //lib/examples:ex\_cheat}. The runner
executes the embedded testbench (Section~\ref{sec:t-testbench}) and emits the
derived Verilog netlist in Listing~\ref{lst:t-counter-out}. Synchronous reset
\code{rst} is bound implicitly. Figure~\ref{fig:t-counter} illustrates the
cycle-accurate trace: \code{n} increments while \code{enable} is asserted,
\code{tick} pulses when \code{n} equals seven, and progression halts when
\code{enable} deasserts.

\begin{figure}[htbp]
\centering
\input{paper_cheat_timing.tex}
\caption{The counter, from its trace: \code{enable} drops at cycle
10, \code{n} counts while it is high, \code{tick} pulses on the
edge after each seven.}
\label{fig:t-counter}
\end{figure}

\lstinputlisting[style=ascii,caption={What \code{ex\_cheat} printed: the counter's Verilog, from the same file.},label={lst:t-counter-out}]{out_cheat.txt}

\section{Step Two: Time}
\label{sec:t-time}

Discrete time is quantified in half-cycle \emph{ticks}. The default clock runs
with a period of two ticks, rising on even ticks and falling on odd ticks.

A \emph{register} updates exclusively at clock edges. When code evaluates
\code{self.n <= n + 1}, the update is scheduled rather than committed
immediately; all registered drives latch simultaneously at the end of the
simulation step. Consequently, reading a register after an edge yields its value
prior to the step, while driven values emerge in the subsequent cycle. Computing
\code{tick} from \code{n} prior to the drive ensures the output reflects the
registered state at the beginning of the cycle, matching hardware RTL behavior.

A \emph{wire} represents combinational nets driven by \code{set}: output ports
or intermediate \code{let} bindings, which lowering preserves as named nets.
Wires hold no state and evaluate strictly as combinational functions of current
inputs and registered state.

The core temporal invariant is simple: \textbf{every \code{.await} defines a
cycle boundary.} A process cannot cross multiple edges in a single simulation
step; a subsequent wait statement suspends until the next active edge.
Evaluating concurrent events at an identical edge requires \code{parallel!}
(simulation only). State sampled prior to a wait reflects the previous cycle;
reads must follow the wait statement.

Multi-step sequential routines execute across successive cycles by placing
multiple wait points within a loop. The lowering engine transforms such
processes into explicit finite state machines, where each wait defines a
distinct state and conditional waits branch across the state graph. Signals
asserted in a state hold their values until reassigned.

\section{Step Three: Choosing}
\label{sec:t-choosing}

Hardware circuits implement multiplexers rather than procedural branches:
competing expressions evaluate concurrently and a selector selects the active
signal. TxHDL provides explicit selection primitives to reflect circuit
topology. Macro \code{with!} specifies conditional register updates, supporting
single guards or prioritized chains with an optional \code{else}. Macro
\code{when!(c => self \{ .. \})} places the condition first for guard-centric
logic. Macro \code{case!} implements multi-way pattern matching, the primary
building block of finite state machines. Standard Rust \code{if} statements
evaluate boolean conditions around \code{set} drives, nesting as priority
decoders. Listing~\ref{lst:t-case} illustrates an explicit state machine
sequencing an execution countdown.

\lstinputlisting[caption={A state machine with \code{case!}: each arm is a Rust pattern, guards included, and the first that matches wins. \code{ex\_case.rs}.},label={lst:t-case}]{lib/examples/ex_case.rs}

Case arms match native Rust patterns in top-to-bottom order, including pattern
guards (\code{State::Run if n == 0}). State representations use custom enums:
deriving \code{Value} automatically establishes netlist bit width (two bits for
four variants) and waveform formatting. Figure~\ref{fig:t-case} demonstrates
cycle progression: assertion of \code{go} triggers transition to \code{Load},
followed by three decrements in \code{Run}, entering \code{Done}, and returning
to \code{Idle}. Listing~\ref{lst:t-case-out} shows execution logs alongside the
emitted Verilog state machine.

\begin{figure}[htbp]
\centering
\input{paper_case_timing.tex}
\caption{The sequencer: \code{go} for a cycle, then \code{Load},
three cycles of \code{Run}, \code{Done}, and back to \code{Idle}.}
\label{fig:t-case}
\end{figure}

\lstinputlisting[style=ascii,caption={What \code{ex\_case} printed: one state per cycle, then the lowering.},label={lst:t-case-out}]{out_case.txt}

For selecting \emph{values} rather than drives, \code{select!} synthesizes an
explicit multiplexer (\code{select!(op => \{ 0 => a + b, 1 => a << b, ..\})}),
as used in the Vreteno ALU. Primitive \code{mux(c, a, b)} provides 2-to-1
selection. Rust value-yielding \code{if/else} blocks and \code{match}
expressions lower to identical multiplexer trees, supporting pattern ranges
(\code{0x30..=0x39}) and range checks (\code{(lo..=hi).contains(\&x)}). In
summary: choose values via \code{select!}, \code{mux}, or \code{match}; choose
drives via \code{with!}, \code{when!}, \code{case!}, or \code{if}.

\section{Step Four: Talking Between Units}
\label{sec:t-channels}

Inter-unit communication relies on \emph{channels}. A channel transports a
structured \emph{transaction} implementing the \code{Transaction} trait,
automatically derived for any \code{U<N>} or composite value struct. Channels
expose paired endpoints: \code{Tx} for sending and \code{Rx} for receiving. The
runtime arbitrates transfers through an elastic 2-element FIFO buffer, ensuring
offered data commits on the subsequent cycle and receivers never sample unready
payloads. Listing~\ref{lst:t-stage} presents an elastic pipeline stage that
synchronizes via \code{until}, reads an incoming transaction, and forwards an
incremented result.

\lstinputlisting[rangebeginprefix=//\ begin\{,rangebeginsuffix=\},rangeendprefix=//\ end\{,rangeendsuffix=\},includerangemarker=false,linerange=unit-unit,caption={A stage between two channels: \code{until} is the wait, and its condition guards everything after it. From \code{ex\_stage.rs}.},label={lst:t-stage}]{lib/examples/ex_stage.rs}

Primitive \code{until(clock, condition)} suspends execution until an active
clock edge on which the boolean predicate evaluates true. Channel methods
include \code{peek} to inspect head data non-destructively, \code{ready} to query
downstream buffer capacity, \code{recv} to dequeue, \code{take} to sample
\code{(valid, data)} opportunistically each cycle, and \code{send} to enqueue.
In synthesised netlists, a channel expands into payload, \code{valid}, and
\code{ready} signals, with the \code{until} predicate synthesizing as an
activation guard. Figure~\ref{fig:t-stage} illustrates the resulting channel
handshake.

\begin{figure}[htbp]
\centering
\input{paper_stage_timing.tex}
\caption{The stage: a word offered on \code{inp} every other cycle,
taken, and passed on incremented on \code{out} a cycle later.}
\label{fig:t-stage}
\end{figure}

Modules running continuously every clock cycle---such as processor
pipelines---avoid blocking \code{until} calls. Instead, they sample channels
opportunistically using \code{take} and drive downstream channels under standard
\code{if} blocks, where the branch condition maps directly to channel
\code{valid}. Listing~\ref{lst:t-tap} shows this pattern in a stream filter
routing odd values.

\lstinputlisting[rangebeginprefix=//\ begin\{,rangebeginsuffix=\},rangeendprefix=//\ end\{,rangeendsuffix=\},includerangemarker=false,linerange=unit-unit,caption={A tap: a receive with no wait, and a send under a plain \code{if}. From \code{ex\_tap.rs}.},label={lst:t-tap}]{lib/examples/ex_tap.rs}

Crucially, raw combinational wires between modules provide no cross-module
evaluation ordering guarantees within the same simulation step. Transactional
channels enforce explicit cycle boundaries. When direct wires are necessary,
drive them strictly from registered outputs and schedule the driver first.

\section{Step Five: Ports With Names}
\label{sec:t-ports}

While scalar ports suffice for simple modules, complex interfaces group signals
within structured structs to prevent accidental pin-swapping bugs. Port bundles
declare channel and wire endpoints as named fields, preserving clarity at
instantiation and extraction. Listing~\ref{lst:t-ports} defines a signed
subtractor whose identical operand types are disambiguated by named fields.

\lstinputlisting[rangebeginprefix=//\ begin\{,rangebeginsuffix=\},rangeendprefix=//\ end\{,rangeendsuffix=\},includerangemarker=false,linerange=unit-unit,caption={Ports as structs: \code{Diff} takes \code{DiffIn} and gives \code{DiffOut}, and \code{Top} takes the sink of an \code{interface!}. From \code{ex\_ports.rs}.},label={lst:t-ports}]{lib/examples/ex_ports.rs}

\code{Diff} accepts a structured \code{DiffIn} bundle and produces
\code{DiffOut}. Function \code{run} destructures input ports in its signature
while addressing output fields by name (\code{out.diff}, \code{out.borrow}).
Port structs support generic parameters and external declarations via
\code{\#[derive(Ports)]}. For example, all AXI-Lite peripherals in
\code{//lib/parts} consume standard \code{LitePort} bundles. Peripheral
registers are declared once via \code{regmap!}, specifying addresses, access
permissions, and bit fields; the macro derives address decoders, C firmware
headers, and datasheet tables automatically.

Port arrays (\code{ins: [Rx<U<8>>; N]}) parameterise module interfaces across
compile-time constants \code{N}. Processes traverse them with unrolled
\code{for i in 0..N} loops, synthesizing indexed netlist ports (\code{ins\_0},
\code{ins\_1}), as seen in the multi-host AXI arbiter.

Macro \code{interface!} declares bus protocols once, synthesizing asymmetric
role structs. \code{Operands} defines a dual-channel link: the \code{Source}
role transmits while \code{Sink} receives, connected via
\code{Operands::new()}. The top-level module maps sink signals to \code{Diff} by
name, making port order irrelevant.

Lowering flattens structured ports into discrete nets.
Listing~\ref{lst:t-ports-out} shows emitted Verilog port declarations with
flattened channel signals (\code{diff\_data}, \code{diff\_valid},
\code{diff\_ready}). Figure~\ref{fig:t-ports} illustrates cycle-accurate
subtraction and borrow generation.

\lstinputlisting[style=ascii,linerange={20-27},caption={The ports of the module \code{ex\_ports} printed: every field, flattened, under its name.},label={lst:t-ports-out}]{out_ports.txt}

\begin{figure}[htbp]
\centering
\resizebox{\columnwidth}{!}{\input{paper_ports_timing.tex}}
\caption{The subtractor: a pair is taken when both operands are
offered and the output has room, and the difference leaves a cycle
later. \code{3 - 7} raises \code{borrow} and leaves 252, \code{fc} in hex.}
\label{fig:t-ports}
\end{figure}

\section{Step Six: Units of Units, and Their Clocks}
\label{sec:t-units}

Modular architectures compose child units within parent structs. A parent holds
child instances as struct fields, wires their ports inside \code{run}, and
schedules concurrent execution using \code{join2} (nested for three or more
units). Lowering generates hierarchical module instantiations. A parent can
forward complex port bundles (\code{bus}) across hierarchy boundaries,
instantiate bidirectional links via \code{link::<LitePort<32, 32, 4>>()}, and tie
unused inputs to constants using \code{tie(v)}.

Because a unit process synchronizes to a single clock, multi-clock systems
mandate hierarchical decomposition. A clock domain is an affine type
implementing \code{Clock}, specifying its name and tick period.
\code{DefaultClock} has a period of two ticks. Listing~\ref{lst:t-twoclock}
defines a clock domain running three times slower alongside two independent
counters. The default counter names its clock only in edge waits. The slow
counter binds \code{Slow} in its wait, registers, and port signatures
(\code{Reg<U<8>, Slow>}, \code{In<Bit, Slow>}). The parent module declares no
waits or internal registers, forwarding clock ports directly. Lowering emits
distinct clock inputs and validates each interface against its domain edges
(Figure~\ref{fig:t-twoclock}).

\lstinputlisting[rangebeginprefix=//\ begin\{,rangebeginsuffix=\},rangeendprefix=//\ end\{,rangeendsuffix=\},includerangemarker=false,linerange={clock-clock,units-units},caption={A clock three times slower, a counter on each clock, and a parent that holds both. From \code{ex\_twoclock.rs}.},label={lst:t-twoclock}]{lib/examples/ex_twoclock.rs}

\begin{figure}[htbp]
\centering
\input{paper_twoclock_timing.tex}
\caption{Two clocks in one module: \code{a} counts on \code{clk} and
\code{b} on \code{slow}, each while its own \code{go} is high.}
\label{fig:t-twoclock}
\end{figure}

Signals cannot cross clock domains directly; type checking enforces clock domain
isolation statically. Inter-domain transfers require explicit dual-clock
synchronizers such as \code{ChanCdc}, implementing dual-port RAM buffers
synchronized by two-stage Gray-code pointer synchronizers, verified under
\code{ex\_cdc}.

Synchronous resets are implicit. Lowered modules receive an active-high
\code{rst} input, asserting which restores initial register values and clears
channel buffers. Modules requiring explicit reset routines declare an input
\code{rst: In<Bit>}, joining the global reset net while executing custom
initialization logic. The Vreteno processor uses this to initialize its program
counter to the reset vector.

\section{Step Seven: Running and Watching}
\label{sec:t-testbench}

The \code{main} function serves as the unit testbench. Listing~\ref{lst:t-main}
shows the counter testbench: instantiating wires via \code{signal}, attaching
waveform logging via \code{Wave::from\_env()}, initializing execution state
with \code{Running::new}, and stepping through simulation cycles via
\code{sim.cycle()}.

\lstinputlisting[rangebeginprefix=//\ begin\{,rangebeginsuffix=\},rangeendprefix=//\ end\{,rangeendsuffix=\},includerangemarker=false,linerange=main-main,caption={The counter's testbench: wires, a waveform, the unit running, twelve cycles. From \code{ex\_cheat.rs}.},label={lst:t-main}]{lib/examples/ex_cheat.rs}

Simulation tracing is controlled by environment variables.
\code{Wave::from\_env()} instantiates a logger when \code{TXHDL\_FST} or
\code{TXHDL\_VCD} is set, allowing binaries to run silently or generate traces
without recompilation. Testbenches register clocks, wire signals, and units
deriving \code{Trace}, logging all active state between \code{start()} and
\code{stop()}. Documentation waveforms are generated via Bazel \code{waveform()}
targets that execute the example and render figures automatically, guaranteeing
consistency with code.

Custom enums and structs deriving \code{Value} render symbolic variant names
directly in traces. Internal combinational nets not bound to ports can be
observed by declaring them as struct fields of type \code{Wire<T>}: driving them
via \code{set} and sampling via \code{get} enables waveform capture while
synthesizing into identical netlist wires.

Units also execute within standard Rust \code{\#[test]} harnesses. The runtime
maintains per-thread executors, executing unit tests concurrently with isolated
simulation clocks. Test assertions should verify registered state following
cycle advancement rather than transient combinational nets.

\section{Step Eight: From Rust to Hardware}
\label{sec:t-netlist}

Listing~\ref{lst:t-counter-out} shows Verilog emitted by
\code{Counter::verilog("counter")}. Equivalent netlist data structures and VHDL
entities are generated via \code{Counter::lowered("counter")}. Both
representations derive from the \code{run} process: \code{\#[lower]} traverses
the AST, mapping registers to clocked flip-flops, \code{let} bindings to named
wires, \code{with!}/\code{case!} to synchronous assignment blocks,
\code{select!}/\code{mux} to multiplexers, edge waits to sensitivity lists, and
\code{Mem<T, N>} to synchronous memory arrays.

Lowerable logic supports standard arithmetic and bitwise operators (\code{+},
\code{-}, \code{*}, \code{\%}, \code{\&}, \code{|}, \code{\^{}}, \code{!},
\code{<<}, \code{>>}), signed comparisons, \code{as} casts, bit extraction
(\code{slice}), concatenation (\code{concat}), sign/zero extension
(\code{sext}/\code{zext}), and multi-limb wide integers (\code{U<256, 2>}). Loops
without internal waits unroll statically; loops containing waits synthesize as
counter-driven state machines (\code{ex\_repeat}). Functions annotated with
\code{\#[lower]} inline transparently, allowing modular helper routines
(\code{ex\_helpers}). Keyword collisions with Verilog or VHDL are resolved by
appending \code{\_rw}.

Generated netlists undergo rigorous co-simulation. For every lowered example,
Bazel captures the golden Rust execution trace, derives VHDL and Verilog
netlists, and generates testbenches that replay inputs and assert register and
port equivalence at every tick under \code{nvc} (VHDL) and Verilator (Verilog).
Custom units inherit these checks automatically upon declaration in
\code{waveform()}.

Beyond regression co-simulation, modules formalize invariants using \code{check!}
(safety assertions), \code{assume!} (environment preconditions), and
\code{cover!} (reachability targets). Simulation halts on violated assertions,
while netlist assertions export to SymbiYosys for formal unbounded proof across
all inputs.

\section{Step Nine: Adding Your Own}
\label{sec:t-adding}

Every reference example in the codebase follows a four-step lifecycle:

\begin{enumerate}
\item \textbf{Implement the Unit:} Create \code{lib/examples/ex\_\emph{name}.rs},
  implementing the module and a testbench that runs simulation and derives
  netlists via \code{txhdl::netlist::write\_vhdl\_from\_env}.
\item \textbf{Register Build Targets:} Declare a \code{rust\_binary} in
  \code{lib/examples/BUILD.bazel} and register the file in the \code{sources}
  filegroup.
\item \textbf{Configure Waveform Verification:} Add a \code{waveform()} target
  in \code{docs/BUILD.bazel} specifying trace signals and declaring
  \code{lowered = (entity, unit)} for automated co-simulation.
\item \textbf{Document the Module:} Author a numbered exposition file in
  \code{docs/examples\_sections/}, registering the generated figure and
  execution log in the examples document.
\end{enumerate}

This workflow enforces the triad rule: every feature exists simultaneously in
the runtime library, in an automated executable test, and in verified
documentation.

\section{Where Next}
\label{sec:t-next}

\code{//docs:examples} catalogues reference designs organized by language
feature. Three specialized tutorials extend this foundation: \code{//docs:zero}
configures an isolated Bazel workspace from scratch; \code{//docs:axi} details
AXI bus interconnects and peripheral routers; \code{//docs:station} examines
out-of-order reservation stations. The theoretical foundations of language
reduction appear in \code{//docs:embedding}~\cite{embedding}; discrete-event
runtime mechanics are documented in \code{//docs:runtime}~\cite{runtime}; the
complete Vreteno RISC-V processor demonstrates these principles scaled to
physical silicon.

\begin{thebibliography}{10}

\bibitem{embedding}
F.~Filmar and D.~Janković, ``Deleting a language: Embedding TxHDL in
Rust,'' project repository \code{HDL/txhdl}, \code{//docs:embedding},
2026.

\bibitem{runtime}
F.~Filmar and D.~Janković, ``The TxHDL runtime,'' project repository
\code{HDL/txhdl}, \code{//docs:runtime}, 2026.

\bibitem{examples}
F.~Filmar and D.~Janković, ``TxHDL by example,'' project repository
\code{HDL/txhdl}, \code{//docs:examples}, 2026.

\bibitem{cover}
F.~Filmar and D.~Janković, ``TxHDL documents,'' project repository
\code{HDL/txhdl}, \code{//docs:cover}, 2026.

\end{thebibliography}

\end{document}
